% 来源及取材范围见分题文件；此为整理者转录的正文，不是下载原件。
\paragraph{Definitions 3.1, 3.2, 3.4, 3.7.}
\noindent\textit{以下背景与记号据所列原文章节摘编；不是逐字引文，也不是新增证明。}\par
Let $X,Y$ be sets of vertices in a graph $G$. Define
\[e_G(X,Y)=|\{(x,y)\in X\times Y:xy\in E(G)\}|,\qquad d_G(X,Y)=\frac{e_G(X,Y)}{|X||Y|}.\]
We drop the subscript when the context is clear. A pair $(X,Y)$ is $\varepsilon$-regular if for all $A\subset X$, $B\subset Y$ with $|A|\geq\varepsilon|X|$, $|B|\geq\varepsilon|Y|$ one has $|d(A,B)-d(X,Y)|\leq\varepsilon$. A partition $\mathcal P=\{V_1,\ldots,V_k\}$ is $\varepsilon$-regular if
\[\sum_{\substack{(i,j)\in[k]^2\\(V_i,V_j)\text{ not }\varepsilon\text{-regular}}}|V_i||V_j|\leq\varepsilon|V(G)|^2.\]
Let $n=|V(G)|$. Put $q(U,W)=|U||W|d(U,W)^2/n^2$. For partitions $\mathcal P_U=\{U_1,\ldots,U_k\}$ and $\mathcal P_W=\{W_1,\ldots,W_l\}$ put
\[q(\mathcal P_U,\mathcal P_W)=\sum_{i=1}^k\sum_{j=1}^lq(U_i,W_j),\qquad q(\mathcal P)=q(\mathcal P,\mathcal P).\]
Observe that $0\leq q(\mathcal P)\leq1$ because the edge density is bounded above by $1$.

\begin{theorem}[3.5, Szemer\'edi's regularity lemma]
For every $\varepsilon>0$, there exists a constant $M$ such that every graph has an $\varepsilon$-regular partition into at most $M$ parts.
\end{theorem}
\begin{lemma}[3.8]
For any partitions $\mathcal P_U$ and $\mathcal P_W$ of vertex sets $U$ and $W$, $q(\mathcal P_U,\mathcal P_W)\geq q(U,W)$.
\end{lemma}
\begin{proof}
Let $\mathcal P_U=\{U_1,\ldots,U_k\}$ and $\mathcal P_W=\{W_1,\ldots,W_l\}$. Choose vertices $x$ uniformly from $U$ and $y$ uniformly from $W$. Let $U_i$ be the part containing $x$ and $W_j$ the part containing $y$. Then define the random variable $Z=d(U_i,W_j)$. Let us look at properties of $Z$. The expectation is
\[\E[Z]=\sum_{i=1}^k\sum_{j=1}^l\frac{|U_i|}{|U|}\frac{|W_j|}{|W|}d(U_i,W_j)=\frac{e(U,W)}{|U||W|}=d(U,W).\]
The second moment is
\[\E[Z^2]=\sum_{i=1}^k\sum_{j=1}^l\frac{|U_i|}{|U|}\frac{|W_j|}{|W|}d(U_i,W_j)^2=\frac{n^2}{|U||W|}q(\mathcal P_U,\mathcal P_W).\]
By convexity, $\E[Z^2]\geq\E[Z]^2$, which implies the lemma.
\end{proof}
\begin{lemma}[3.9]
If $\mathcal P'$ refines $\mathcal P$, then $q(\mathcal P')\geq q(\mathcal P)$.
\end{lemma}
\begin{proof}
Let $\mathcal P=\{V_1,\ldots,V_m\}$ and apply Lemma 3.8 to every $(V_i,V_j)$.
\end{proof}
\begin{lemma}[3.10, Energy boost lemma]
If $(U,W)$ is not $\varepsilon$-regular as witnessed by $U_1\subset U$ and $W_1\subset W$, then
\[q(\{U_1,U\setminus U_1\},\{W_1,W\setminus W_1\})>q(U,W)+\varepsilon^4\frac{|U||W|}{n^2}.\]
\end{lemma}
\begin{proof}
Define $Z$ as in the proof of Lemma 3.8. Then
\[\Var(Z)=\E[Z^2]-\E[Z]^2=\frac{n^2}{|U||W|}\big(q(\{U_1,U\setminus U_1\},\{W_1,W\setminus W_1\})-q(U,W)\big).\]
But observe that $|Z-\E[Z]|=|d(U_1,W_1)-d(U,W)|$ with probability $|U_1||W_1|/(|U||W|)$ (corresponding to $x\in U_1$ and $y\in W_1$), so
\[\begin{split}
\Var(Z)=\E[(Z-\E[Z])^2]&\geq\frac{|U_1|}{|U|}\frac{|W_1|}{|W|}(d(U_1,W_1)-d(U,W))^2\\
&>\varepsilon\cdot\varepsilon\cdot\varepsilon^2,
\end{split}\]
as desired.
\end{proof}
\begin{lemma}[3.11]
If a partition $\mathcal P=\{V_1,\ldots,V_k\}$ of $V(G)$ is not $\varepsilon$-regular, then there exists a refinement $\mathcal Q$ of $\mathcal P$ where every $V_i$ is partitioned into at most $2^k$ parts such that $q(\mathcal Q)\geq q(\mathcal P)+\varepsilon^5$.
\end{lemma}
\begin{proof}
For all $(i,j)$ such that $(V_i,V_j)$ is not $\varepsilon$-regular, find $A_{i,j}\subset V_i$ and $A_{j,i}\subset V_j$ that witness irregularity (do this simultaneously for all irregular pairs). Let $\mathcal Q$ be a common refinement of $\mathcal P$ by $A_{i,j}$'s. Each $V_i$ is partitioned into at most $2^k$ parts as desired. Then
\[\begin{split}
q(\mathcal Q)&=\sum_{(i,j)\in[k]^2}q(\mathcal Q_{V_i},\mathcal Q_{V_j})\\
&=\sum_{\substack{(i,j)\in[k]^2\\(V_i,V_j)\ \varepsilon\text{-regular}}}q(\mathcal Q_{V_i},\mathcal Q_{V_j})
+\sum_{\substack{(i,j)\in[k]^2\\(V_i,V_j)\text{ not }\varepsilon\text{-regular}}}q(\mathcal Q_{V_i},\mathcal Q_{V_j}),
\end{split}\]
where $\mathcal Q_{V_i}$ is the partition of $V_i$ given by $\mathcal Q$. By Lemma 3.8, the above quantity is at least
\[\begin{split}
&\sum_{\substack{(i,j)\in[k]^2\\(V_i,V_j)\ \varepsilon\text{-regular}}}q(V_i,V_j)\\
&\quad+\sum_{\substack{(i,j)\in[k]^2\\(V_i,V_j)\text{ not }\varepsilon\text{-regular}}}
q(\{A_{i,j},V_i\setminus A_{i,j}\},\{A_{j,i},V_j\setminus A_{j,i}\}),
\end{split}\]
since $V_i$ is cut by $A_{i,j}$ when creating $\mathcal Q$, so $\mathcal Q_{V_i}$ is a refinement of $\{A_{i,j},V_i\setminus A_{i,j}\}$. By Lemma 3.10, the above sum is at least
\[\sum_{(i,j)\in[k]^2}q(V_i,V_j)+\sum_{\substack{(i,j)\in[k]^2\\(V_i,V_j)\text{ not }\varepsilon\text{-regular}}}\varepsilon^4\frac{|V_i||V_j|}{n^2}.\]
But the second sum is at least $\varepsilon^5$ since $\mathcal P$ is not $\varepsilon$-regular, so we deduce the desired inequality.
\end{proof}
\begin{proof}[Proof of Theorem 3.5]
Start with a trivial partition. Repeatedly apply Lemma 3.11 whenever the current partition is not $\varepsilon$-regular. By the definition of energy, $0\leq q(\mathcal P)\leq1$. However, by Lemma 3.11, $q(\mathcal P)$ increases by at least $\varepsilon^5$ at each iteration. So we will stop after at most $\varepsilon^{-5}$ steps, resulting in an $\varepsilon$-regular partition.
\end{proof}

\begin{theorem}[3.13, Triangle counting lemma]
Let $G$ be a graph and $X,Y,Z$ be subsets of the vertices of $G$ such that $(X,Y),(Y,Z),(Z,X)$ are all $\varepsilon$-regular pairs some $\varepsilon>0$. Let $d_{XY},d_{XZ},d_{YZ}$ denote the edge densities $d(X,Y),d(X,Z),d(Y,Z)$ respectively. If $d_{XY},d_{XZ},d_{YZ}\geq2\varepsilon$, then the number of triples $(x,y,z)\in X\times Y\times Z$ such that $x,y,z$ form a triangle in $G$ is at least
\[(1-2\varepsilon)(d_{XY}-\varepsilon)(d_{XZ}-\varepsilon)(d_{YZ}-\varepsilon)|X||Y||Z|.\]
\end{theorem}
\begin{proof}
By assumption, $(X,Y)$ is an $\varepsilon$-regular pair. This implies that fewer than $\varepsilon|X|$ of the vertices in $X$ have fewer than $(d_{XY}-\varepsilon)|Y|$ neighbors in $Y$. If this were not the case, then we could take $Y$ together with the subset consisting of all vertices in $X$ that have fewer than $(d_{XY}-\varepsilon)|Y|$ neighbors in $Y$ and obtain a pair of subsets witnessing the irregularity of $(X,Y)$, which would contradict our assumption. Intuitively these bounds make sense, since if the edges between $X$ and $Y$ were random-like we would expect most vertices in $X$ to have about $d_{XY}|Y|$ neighbors in $Y$, meaning that not too many vertices in $X$ can have very small degree in $Y$.

Applying the same argument to the $\varepsilon$-regular pair $(X,Z)$ proves the analogous result that fewer than $\varepsilon|X|$ of the vertices in $X$ have fewer than $(d_{XZ}-\varepsilon)|Z|$ neighbors in $Z$. Combining these two results, we see that we can find a subset $X'$ of $X$ of size at least $(1-2\varepsilon)|X|$ such that every vertex $x\in X'$ is adjacent to at least $(d_{XY}-\varepsilon)|Y|$ of the elements in $Y$ and $(d_{XZ}-\varepsilon)|Z|$ of the elements in $Z$. Using the hypothesis that $d_{XY},d_{XZ}\geq2\varepsilon$ and the fact that $(Y,Z)$ is $\varepsilon$-regular, we see that for any $x\in X'$, the edge density between the neighborhoods of $x$ in $Y$ and $Z$ is at least $(d_{YZ}-\varepsilon)$.

Now, for each vertex $x\in X'$, of which there are at least $(1-2\varepsilon)|X|$, and choice of edge between the neighborhoods of $x$ in $Y$ and $x$ in $Z$, of which there are at least $(d_{XY}-\varepsilon)(d_{XZ}-\varepsilon)(d_{YZ}-\varepsilon)|Y||Z|$, we get a unique $(X,Y,Z)$-triangle in $G$. It follows that the number of such triangles is at least
\[(1-2\varepsilon)(d_{XY}-\varepsilon)(d_{XZ}-\varepsilon)(d_{YZ}-\varepsilon)|X||Y||Z|\]
as claimed.
\end{proof}

\begin{theorem}[3.15, Triangle removal lemma]
For all $\varepsilon>0$, there exists $\delta>0$ such that any graph on $n$ vertices with less than or equal to $\delta n^3$ triangles can be made triangle-free by removing at most $\varepsilon n^2$ edges.
\end{theorem}
\begin{proof}
Suppose we are given a graph on $n$ vertices with fewer than $\delta n^3$ triangles, for some parameter $\delta$ we will choose later. Begin by taking an $\varepsilon/4$-regular partition of the graph with parts $V_1,V_2,\ldots,V_M$. Next, for each ordered pair of parts $(V_i,V_j)$, remove all edges between $V_i$ and $V_j$ if
\begin{enumerate}[label=(\alph*)]
\item $(V_i,V_j)$ is an irregular pair,
\item the density $d(V_i,V_j)$ is less than $\varepsilon/2$, or
\item either $V_i$ or $V_j$ has at most $(\varepsilon/4M)n$ vertices (is ``small'').
\end{enumerate}
How many edges are removed in this process? Well, since we took an $\varepsilon/4$-regular partition, by definition
\[\sum_{\substack{i,j\\(V_i,V_j)\text{ not }(\varepsilon/4)\text{-regular}}}|V_i||V_j|\leq\frac{\varepsilon}{4}n^2.\]
So at most $(\varepsilon/4)n^2$ edges are removed between irregular pairs in (a). The number of edges removed from low-density pairs in (b) is
\[\sum_{\substack{i,j\\d(V_i,V_j)<\varepsilon/2}}d(V_i,V_j)|V_i||V_j|
\leq\frac{\varepsilon}{2}\sum_{i,j}|V_i||V_j|=\frac{\varepsilon}{2}n^2,\]
where the intermediate sum is taken over all ordered pairs of parts. The number of edges removed between small parts in (c) is at most
\[n\cdot\frac{\varepsilon}{4M}n\cdot M=\frac{\varepsilon}{4}n^2,\]
since each of the $n$ vertices is adjacent to at most $(\varepsilon/4M)n$ vertices in each small part, and there are at most $M$ small parts.

As expected, cleaning up the graph by removing edges between badly behaving parts does not remove too many edges. We claim that after this process, for some choice of $\delta$, the graph is triangle-free. The removal lemma follows from this claim, since the previous step removed less than $\varepsilon n^2$ edges from the graph.

Indeed, suppose that after following the above procedure and (possibly) removing some edges the resulting graph still has some triangle. Then we can find parts $V_i,V_j,V_k$ (not necessarily distinct) containing each of the vertices of this triangle. Because edges between the pairs described in (a) and (b) were removed, $V_i,V_j,V_k$ satisfy the hypotheses of the triangle counting lemma. Applying Theorem 3.13 to this triple of subsets implies that the graph still has at least
\[\left(1-\frac{\varepsilon}{2}\right)\left(\frac{\varepsilon}{4}\right)^3|V_i||V_j||V_k|\]
such triangles. By (c) each of these parts has size at least $(\varepsilon/4M)n$, so in fact the number of $(V_i,V_j,V_k)$-triangles after removal is at least
\[\left(1-\frac{\varepsilon}{2}\right)\left(\frac{\varepsilon}{4}\right)^3\left(\frac{\varepsilon}{4M}\right)^3n^3.\]
Then by choosing positive
\[\delta<\frac16\left(1-\frac{\varepsilon}{2}\right)\left(\frac{\varepsilon}{4}\right)^3\left(\frac{\varepsilon}{4M}\right)^3\]
we obtain a contradiction, since the original graph has less than $\delta n^3$ triangles by assumption, but the triangle counting lemma shows that we have strictly more than this many triangles after removing some edges in the graph. The factor of $1/6$ is included here to deal with overcounting that may occur (e.g. when $V_i=V_j=V_k$). Since $\delta$ only depends on $\varepsilon$ and the constant $M$ from Theorem 3.5, this completes our proof.
\end{proof}
